\section{Motivation und leitende Unterscheidung}

Die übliche Darstellung einer Ellipse lautet
\[
  x(\theta)=c_x+a\cos\theta,
  \qquad
  y(\theta)=c_y+b\sin\theta.
\]
Diese Schreibweise legt nahe, die trigonometrischen Funktionen seien ein
wesentlicher Bestandteil der Ellipse. Tatsächlich legen sie aber zugleich zwei
verschiedene Dinge fest: die geometrische Punktmenge und die Geschwindigkeit,
mit der diese Punktmenge durchlaufen wird.

Das hier untersuchte Modell trennt beide Ebenen. Eine skalare Steuerfunktion
liefert nur die horizontale Lage. Die vertikale Lage wird anschließend so
rekonstruiert, dass die Ellipsengleichung invariant bleibt. Damit können eine
lineare Dreieckswelle, geglättete Zwischenformen und der Kosinus dieselbe
geometrische Kurve erzeugen, obwohl sie diese verschieden parametrisieren.

Die leitende These lautet daher:
\begin{quote}
Die Glättung verändert nicht die Ellipse, sondern ihre innere Uhr.
\end{quote}
